\section{Step 2: completing \texorpdfstring{$T$}{T} away from the ports}\label{sec:completion}
Step~(ii) of Eq.~\eqref{eq:two-steps} fixes the fitted
$\theta=(W,\omega,\Gamma)$ and selects $T$ from the family
\begin{equation}
 \mathcal T(\theta)=\{T\in\R^{n\times r}:AT=W\}.
 \label{eq:family}
\end{equation}
By Eq.~\eqref{eq:ports-see} every member reproduces the fitted FRFs. Two
requirements select the member. It must keep every modal coordinate, and
the lifted motion $u=Tq$ should obey the nominal equation of motion.

\subsection{Keeping every modal coordinate}
Port matching alone need not give $T$ independent columns, because
$W\in\R^{h\times r}$ can have fewer rows than columns. We therefore require
\begin{equation}
 \sigma_r(RT)\ge\delta>0,
 \label{eq:margin}
\end{equation}
where $R\in\R^{n_{\rm d}\times n}$ selects the $n_{\rm d}\ge r$ represented
translation rows of $u$. Then $RT$ has full column rank and
\begin{equation}
 V=(RT)^{\dagger}R,\qquad VT=I_r,\qquad\|V\|_2\le\delta^{-1}.
 \label{eq:completion-left-inverse}
\end{equation}
For a field $u=Tq$, $V$ recovers $q=Vu$ from the represented translations
by least squares, with error amplification at most $1/\delta$.

\subsection{Loss function and constraints}
The loss measures how far the lifted motion $u=Tq$ is from obeying the
nominal equation of motion. Under the measured forces the fitted model
moves as $q(s)=\Delta_\theta(s)^{-1}W^{\top}R_{\rm i}^{\top}f(s)$
(Eq.~\eqref{eq:model}). Insert $u=Tq$ into the nominal equation with the
fitted damping force $M_0T\Gamma\dot q$ and the load $M_0TT^{\top}Bf$ that
the $r$ modes carry. By Eq.~\eqref{eq:latent} the residual force is
\begin{equation}
 \begin{aligned}
 M_0\ddot u+M_0T\Gamma\dot q+K_0u-M_0TT^{\top}Bf
 &=M_0T\bigl(\ddot q+\Gamma\dot q+\Omega^2q-T^{\top}Bf\bigr)+Gq=Gq,\\
 G&=K_0T-M_0T\Omega^2.
 \end{aligned}
 \label{eq:dynamic-defect}
\end{equation}
The loss is the strain energy of this residual relative to the strain
energy of the fitted model. Both are summed over the frequency lines
$\nu_\ell$ of the test, with $q_\ell=q(s_\ell)$:
\begin{equation}
 J_F(\theta,T)
 =\frac{\sum_\ell\|Gq_\ell\|_{K_0^{-1}}^2}{\sum_\ell q_\ell^*\Omega^2q_\ell}
 =\frac{\tr\bigl(G^{\top}K_0^{-1}G\,\Phi_q\bigr)}{\tr(\Omega^2\Phi_q)}.
 \label{eq:dynamic-objective}
\end{equation}
Here $\|x\|_X^2=x^*Xx$. A force $g$ causes the static deflection
$K_0^{-1}g$, and $\|g\|_{K_0^{-1}}^2$ is twice its strain energy. The second
form holds because $G$ does not depend on $\ell$. Both sums then need only
the $r\times r$ modal response matrix
\begin{equation}
 \Phi_q=\operatorname{Re}\sum_\ell\Delta_\theta(s_\ell)^{-1}W^{\top}R_{\rm i}^{\top}
 \Phi_{ff}(\nu_\ell)R_{\rm i}W\Delta_\theta(s_\ell)^{-*}\,\Delta\nu\succeq0.
 \label{eq:Pi}
\end{equation}
It is $\operatorname{Re}\sum_\ell q_\ell q_\ell^*\Delta\nu$ with the force
spectral matrices $\Phi_{ff}(\nu_\ell)$ of the training record in place of
$f(s_\ell)f(s_\ell)^*$. Here $\Delta\nu$ is the line spacing. The matrix
$\Phi_q$ records how much the test moves each mode. It brings the
identified damping and the measured forces into the loss. It does not
depend on $T$, so $J_F$ is quadratic in $T$. The loss is dimensionless and
vanishes when
\begin{equation}
 K_0T=M_0T\Omega^2,
 \label{eq:intertwining}
\end{equation}
that is, when every column of $T$ is a nominal mode shape at its fitted
frequency. An imperfect nominal model generally admits no such $T$ with
$AT=W$. So the loss is minimized subject to three constraints:
\begin{equation}
 \boxed{\ \min_{T}\ J_F(\theta,T)\quad\text{subject to}\quad
 AT=W,\quad\sigma_r(RT)\ge\delta,\quad t_1,\ldots,t_r\in\mathcal U.\ }
 \label{eq:completion-problem}
\end{equation}
The port constraint keeps the fitted FRFs. The margin keeps every modal
coordinate. The columns $t_j$ of $T$ lie in a trial space $\mathcal U$ with
$A\mathcal U=\R^h$, so that every port value is reachable. The chain uses
$\mathcal U=\R^3$. The laboratory space holds the static deflections under
the 33 unit port forces and 384 compliant fields that no port sees
(Appendix~\ref{sec:field-selection}).

\subsection{Feasibility and optimality}

The trial space must leave room for rank $r$. A dependence $Tv=0$ forces
$Wv=ATv=0$, so only $\ker W$ is at risk, and there $T$ maps into
$\mathcal U\cap\ker A$, of dimension $\dim\mathcal U-h$. A member of rank
$r$ therefore exists if and only if
\begin{equation}
 \dim\mathcal U-h\ \ge\ r-\rank W.
 \label{eq:rank-feasibility}
\end{equation}
Necessity is this injectivity. For sufficiency take any $T_1$ with
$AT_1=W$, keep it on the orthogonal complement of $\ker W$ and map
$\ker W$ injectively into $\mathcal U\cap\ker A$. The chain has $1\ge1$
and the laboratory $384\ge351$. The count guarantees rank $r$, not the
margin, which is checked on the result.

Write
\[
 \langle X,Y\rangle_{\Phi_q}=\tr(X^{\top}K_0^{-1}Y\Phi_q),\qquad
 \mathcal L(T)=K_0T-M_0T\Omega^2,\qquad
 J_F=\frac{\langle\mathcal L(T),\mathcal L(T)\rangle_{\Phi_q}}{\tr(\Omega^2\Phi_q)},
\]
a convex quadratic in $T$. Let $\mathcal D=\{D:AD=0,\ d_j\in\mathcal U\}$ be
the directions that keep the ports.

\begin{theorem}[Completion]\label{thm:completion}
Let $M_0,K_0\succ0$, $\Phi_q\succeq0$, $\omega_j>0$ and $\delta>0$.
\begin{enumerate}
\item[(i)] Every feasible $T$ of Eq.~\eqref{eq:completion-problem} has rank
$r$ and satisfies Eq.~\eqref{eq:completion-left-inverse}.
\item[(ii)] Without the margin a minimizer $T_{\rm LS}$ exists. A member $T$
of $\mathcal T(\theta)$ with columns in $\mathcal U$ is one if and only if
$\langle\mathcal L(T),\mathcal L(D)\rangle_{\Phi_q}=0$ for every
$D\in\mathcal D$.
\item[(iii)] Every such member satisfies
$J_F(\theta,T)=J_F(\theta,T_{\rm LS})
+\langle\mathcal L(T-T_{\rm LS}),\mathcal L(T-T_{\rm LS})\rangle_{\Phi_q}/\tr(\Omega^2\Phi_q)$.
\end{enumerate}
\end{theorem}

\begin{proof}
(i) The margin gives $\|RTv\|_2\ge\delta\|v\|_2$, so $RT$ has full column
rank, $(RT)^{\dagger}RT=I_r$, and $\|V\|_2=1/\sigma_r(RT)\le1/\delta$
because the rows of $R$ are orthonormal.

(ii) The admissible set is the nonempty affine set $T_1+\mathcal D$. On it
$J_F$ is a convex quadratic bounded below by zero, so a minimizer exists.
A convex quadratic is minimal at $T$ exactly when its directional
derivative $2\langle\mathcal L(T),\mathcal L(D)\rangle_{\Phi_q}/\tr(\Omega^2\Phi_q)$
vanishes for every admissible direction $D$.

(iii) $T-T_{\rm LS}\in\mathcal D$ and $\mathcal L$ is linear. Expanding
$\mathcal L(T)=\mathcal L(T_{\rm LS})+\mathcal L(T-T_{\rm LS})$, the cross
term vanishes by (ii).
\end{proof}

So $T_{\rm LS}$ solves Eq.~\eqref{eq:completion-problem} whenever it meets
the margin. Otherwise, by (iii), the problem asks for the independent map
closest to $T_{\rm LS}$ in the seminorm of $\mathcal L$. The margin makes it
nonconvex. We start from $T_{\rm LS}$ with its small singular values raised
to the margin and run an alternating-direction iteration that keeps $AT=W$
in every iterate (Appendix~\ref{sec:constrained}). The accepted map is the
best feasible iterate. Its value bounds the optimum,
$J_F(\theta,T_{\rm LS})\le J_F^\star\le J_F(\theta,T)$, so the gap to
$J_F(\theta,T_{\rm LS})$ certifies how much the margin costs.
